\TFMChapter{Introducción}

\TFMSection{Planteamiento del Problema}
La creación de propuestas de viaje personalizadas requiere consultar numerosas fuentes: documentos de proveedores, catálogos de actividades, tarifas, alojamientos, vuelos y datos de cada destino. En muchas agencias esta información está distribuida en PDFs, correos, repositorios internos y plataformas externas, dificultando su localización y reutilización, y el personal debe recopilar las preferencias del cliente, buscar actividades, consultar precios y disponibilidad, y organizarlo todo en una propuesta coherente --un proceso con una gran carga de trabajo manual.

Los sistemas de búsqueda tradicionales se basan en coincidencias exactas de palabras, por lo que pueden no recuperar información relevante si la consulta usa términos diferentes a los empleados; además, los precios de vuelos y alojamientos son datos dinámicos que pueden no estar en la documentación interna o quedar desactualizados (Capítulo 2, ``Recuperación aumentada por generación'').

\TFMSection{Propuesta y Motivación}
En este contexto, las arquitecturas RAG permiten combinar modelos de lenguaje con una base documental propia \parencite{lewis2020rag}: antes de responder, el sistema recupera la información más relevante y la incorpora al contexto del modelo, permitiendo consultar la documentación de proveedores en lenguaje natural con respuestas fundamentadas (fundamentos en el Capítulo 2).

El proyecto plantea además agentes de IA especializados que coordinan tareas --recuperar información turística, consultar precios, organizar el itinerario, generar la imagen de portada y elaborar la propuesta comercial--, basándose prioritariamente en información corporativa y recurriendo a fuentes externas solo cuando sea necesario (diseño detallado en el Capítulo 3).

A partir de esta necesidad, el presente Trabajo Fin de Máster propone el desarrollo de una prueba de concepto destinada a asistir a los profesionales de una agencia de viajes en la creación de propuestas personalizadas. La solución combina procesamiento documental, recuperación de información, modelos de lenguaje y una arquitectura basada en agentes.

La principal motivación es reducir el tiempo dedicado a tareas repetitivas de búsqueda, consulta y organización de información. El sistema no pretende sustituir al profesional, sino proporcionarle una herramienta de apoyo que genere una primera propuesta revisable y facilite el acceso al conocimiento disponible.

Desde el punto de vista empresarial, la solución puede mejorar la eficiencia del proceso comercial, favorecer la reutilización de la documentación de proveedores y reducir la dependencia del conocimiento individual de cada empleado. Desde el punto de vista tecnológico, permite estudiar la aplicación conjunta de sistemas RAG, modelos generativos y agentes inteligentes dentro de una solución orientada a un caso de uso real.

\TFMSection{Objetivos Generales y Específicos}
El objetivo principal ha sido el diseño y desarrollo de un sistema de IA generativa capaz de asistir a los profesionales de una agencia de viajes en la elaboración de propuestas personalizadas, procesando documentación turística, recuperando información mediante RAG y coordinando agentes especializados para destinos, actividades, alojamientos, vuelos y precios, integrando modelos de lenguaje, servicios cloud de AWS y fuentes externas.

Para poder cumplir con el objetivo principal, se han definido los siguientes objetivos específicos:
\begin{itemize}
\item \textbf{Desarrollo de una interfaz básica en Streamlit:} Implementar una interfaz sencilla y funcional que permita introducir las características del viaje, interactuar con el sistema y visualizar los resultados generados como parte de la prueba de concepto.
\item \textbf{Implementación de un pipeline serverless en AWS de procesamiento documental:} Automatizar la ingesta y transformación de documentos turísticos, convirtiendo su contenido en información estructurada y preparada para su explotación.
\item \textbf{Extracción y normalización de información:} Utilizar modelos de inteligencia artificial generativa para extraer texto, metadatos, actividades, destinos y precios de documentos con diferentes formatos y estructuras.
\item \textbf{Diseño e implementación de un sistema RAG:} Desarrollar un mecanismo de recuperación semántica que permita localizar información relevante dentro de la documentación interna y utilizarla como contexto para generar respuestas con valor.
\item \textbf{Desarrollo de una arquitectura multiagente:} Implementar agentes especializados capaces de recopilar los requisitos del usuario, consultar información turística, obtener precios y elaborar una propuesta de viaje coherente.
\end{itemize}


\begin{itemize}
\item \textbf{Integración con fuentes externas:} Incorporar servicios externos, en particular mediante SearchAPI, para consultar información dinámica sobre vuelos y alojamientos cuando los datos disponibles en la base documental no sean suficientes.
\item \textbf{Desarrollo de una API de integración:} Implementar los servicios necesarios para comunicar la interfaz de usuario, el sistema multiagente, la base de conocimiento y los servicios externos.
\item \textbf{Generación de contenido visual para la propuesta:} Incorporar la generación automática de una imagen de portada relacionada con el destino y las actividades seleccionadas, con el fin de enriquecer visualmente la propuesta de viaje.
\item \textbf{Despliegue de la solución en infraestructura cloud:} Configurar los servicios de almacenamiento, procesamiento, indexación y persistencia necesarios para ejecutar la solución en AWS.
\item \textbf{Pruebas y validación del sistema:} Evaluar el funcionamiento del procesamiento documental, la recuperación de información, la coordinación de los agentes y la calidad de las propuestas generadas.
\end{itemize}

El cumplimiento de estos objetivos se verifica en el Capítulo 6 (Pruebas y Validación del Producto).

\TFMSection{Estructura del Proyecto}
La memoria se organiza en ocho capítulos, cuyo contenido se resume en la Tabla \ref{tab:estructura}.

\begin{table}[H]
\centering
\begin{tabular}{@{}p{4.5cm}p{9cm}@{}}
\toprule
\textbf{Capítulo} & \textbf{Contenido} \\
\midrule
1. Introducción & Contexto, motivación, objetivos y alcance del proyecto. \\
2. Estado del Arte y Fundamentos Teóricos & Fundamentos teóricos (modelos de lenguaje, RAG, agentes) y revisión de soluciones existentes. \\
3. Metodología, Requisitos y Diseño de la Solución & Metodología de trabajo, análisis de requisitos y diseño de la arquitectura. \\
4. Implementación & Desarrollo técnico de los componentes del sistema. \\
5. Evaluación y Experimentos Realizados & Evaluación de la solución mediante experimentos. \\
6. Pruebas y Validación del Producto & Pruebas funcionales y validación del producto. \\
7. Discusión & Análisis de los resultados obtenidos y de las limitaciones del proyecto. \\
8. Conclusiones y Trabajo Futuro & Conclusiones y líneas de trabajo futuro. \\
\bottomrule
\end{tabular}
\caption{Estructura de la memoria por capítulos.}
\label{tab:estructura}
\end{table}


\TFMSection{Alcance de la Solución}
El proyecto comprende una prueba de concepto capaz de procesar documentación turística, extraer y normalizar su contenido, indexarlo en una base vectorial y recuperarlo mediante RAG; sobre esta base, un sistema de agentes recopila requisitos del viaje, consulta destinos y actividades, obtiene precios cuando es necesario, genera una imagen de portada y compone la propuesta. Incluye una interfaz básica en Streamlit, una API de comunicación entre componentes y una infraestructura cloud en AWS para almacenamiento, procesamiento, indexación y persistencia.

Al tratarse de una prueba de concepto, quedan fuera del alcance la realización de reservas o pagos, la integración completa con sistema comerciales en producción, la garantía de disponibilidad en tiempo real y el tratamiento exhaustivo de todos los posibles proveedores y destino. Las propuestas generadas deben considerarse una primera versión sujeta a revisión por parte de un profesional.
